\FloatBarrier
\section{Discussion and conclusion}
\label{sec:discussion}

The SDG results separate a uniform computational regime from an exact
pointwise characterization.  If \(\beta_x>\rho\), the primal comparison
problem is strongly convex, the self-centered SDG vanishes exactly at KKT
pairs, and its value controls feasibility and proximal displacement.  At an
arbitrary positive curvature, zero gap is equivalent to global quadratic
support of the fixed Lagrangian, measured by \(\beta_{\crit}(x,y)\).  The
threshold can vary across KKT branches and explains the false numerical zeros
produced by a nonglobal local inner solution.  The stronger zero
interpretation of the global augmented-Lagrangian gap requires a global
nonconvex value problem.

For post-KKT assessment, proximal escape uses a second safe construction: the
hard-constrained proximal map of \(f+\ind_{\Cset}\), distinct from the
Lagrangian proximal comparison in the SDG.  Its ideal radius is infinite at a
global minimizer, finite and positive at a nonglobal local minimizer, and zero
at a point that is not locally minimizing.  Complete stationary catalogs
assess branch recovery, the separable box problem verifies the radius
trichotomy, and the Hessian-flat illustration shows that realized value decrease
can reject stationary points with vanishing gradient and Hessian.

The conclusions assume global \(\rho\)-weak convexity, affine constraints, and
reliable evaluation of the relevant safe proximal subproblems.  Computing
\(\beta_{\crit}(x,y)\) remains a global problem.  The finite rule in
\Cref{alg:safe-finite-escape} accepts a verified lower feasible value as a
rejection certificate and records every other finite-search outcome as
unresolved.
